\documentclass[12pt,a4paper]{article}
\usepackage[utf8]{inputenc}
\usepackage[T1]{fontenc}
\usepackage[french]{babel}
\usepackage{geometry}
\usepackage{xcolor}
\usepackage{fancyhdr}

\geometry{margin=2cm}
\pagestyle{fancy}
\fancyhf{}
\fancyhead[L]{CEJM - Fiche de Révision}
\fancyhead[R]{Chapitre 2}
\fancyfoot[C]{\thepage}
\setlength{\headheight}{15pt}

% Couleurs
\definecolor{bleu}{RGB}{0,90,170}
\definecolor{vert}{RGB}{0,150,0}
\definecolor{rouge}{RGB}{200,0,0}
\definecolor{gris}{RGB}{240,240,240}

% Commandes pour encadrés
\newcommand{\definition}[2]{\noindent\fcolorbox{bleu}{gris}{\parbox{\dimexpr\textwidth-2\fboxsep-2\fboxrule\relax}{\textbf{\textcolor{bleu}{DÉFINITION :}} #1\\\textit{#2}}}}

\newcommand{\pointcle}[1]{\noindent\fcolorbox{vert}{white}{\parbox{\dimexpr\textwidth-2\fboxsep-2\fboxrule\relax}{\textbf{\textcolor{vert}{POINT CLÉ :}} #1}}}

\newcommand{\exemple}[1]{\noindent\fcolorbox{rouge}{white}{\parbox{\dimexpr\textwidth-2\fboxsep-2\fboxrule\relax}{\textbf{\textcolor{rouge}{EXEMPLE :}} #1}}}

\title{\textbf{Fiche de Révision CEJM}\\[2mm]\large Chapitre 2 : Le fonctionnement des marchés}
\author{}
\date{}

\begin{document}
\maketitle

\section*{Problématique}
\textbf{Comment fonctionnent les mécanismes de marché et comment se forme le prix d'équilibre ?}

\section{Définitions essentielles}

\definition{Marché}{Lieu de rencontre entre l'offre (vendeurs) et la demande (acheteurs) où se détermine le prix d'un bien ou service.}

\vspace{3mm}
\definition{Offre}{Quantité de biens ou services que les producteurs sont disposés à vendre à un prix donné durant une période donnée.}

\vspace{3mm}
\definition{Demande}{Quantité de biens ou services que les consommateurs sont disposés à acheter à un prix donné durant une période donnée.}

\vspace{3mm}
\definition{Prix d'équilibre}{Prix auquel l'offre égale la demande. Il n'y a ni pénurie ni excédent.}

\vspace{3mm}
\definition{Élasticité-prix}{Mesure la sensibilité de la demande aux variations de prix. Si élastique, une petite variation de prix entraîne une forte variation de demande.}

\section{Les lois du marché}

\subsection{Loi de l'offre}
\pointcle{Plus le prix augmente, plus l'offre augmente (relation positive)}
\begin{itemize}
    \item Les producteurs sont incités à produire plus quand les prix sont élevés
    \item Recherche de profit maximum
    \item Courbe d'offre croissante
\end{itemize}

\subsection{Loi de la demande}
\pointcle{Plus le prix augmente, plus la demande diminue (relation négative)}
\begin{itemize}
    \item Les consommateurs achètent moins quand les prix sont élevés
    \item Contrainte budgétaire des ménages
    \item Courbe de demande décroissante
\end{itemize}

\subsection{Équilibre du marché}
\pointcle{L'équilibre se forme à l'intersection de l'offre et de la demande}
\begin{itemize}
    \item Prix d'équilibre : aucune tension sur le marché
    \item Quantité d'équilibre : toute l'offre trouve un acheteur
    \item Autorégulation du marché par le mécanisme des prix
\end{itemize}

\section{Les déséquilibres de marché}

\subsection{Situation d'excédent}
\pointcle{Offre > Demande : les prix ont tendance à baisser}
\begin{itemize}
    \item Prix trop élevé par rapport à l'équilibre
    \item Stocks invendus chez les producteurs
    \item Pression à la baisse des prix
\end{itemize}

\subsection{Situation de pénurie}
\pointcle{Demande > Offre : les prix ont tendance à augmenter}
\begin{itemize}
    \item Prix trop bas par rapport à l'équilibre
    \item Demande insatisfaite des consommateurs
    \item Pression à la hausse des prix
\end{itemize}

\section{Les facteurs influençant l'offre et la demande}

\subsection{Facteurs influençant la demande}
\begin{itemize}
    \item \textbf{Revenus des consommateurs} : hausse des revenus → hausse de la demande
    \item \textbf{Prix des biens substituts} : hausse du prix d'un substitut → hausse de la demande
    \item \textbf{Prix des biens complémentaires} : hausse du prix d'un complémentaire → baisse de la demande
    \item \textbf{Goûts et préférences} : effet de mode, publicité
    \item \textbf{Anticipations} : anticipation de hausse → achat immédiat
\end{itemize}

\subsection{Facteurs influençant l'offre}
\begin{itemize}
    \item \textbf{Coûts de production} : baisse des coûts → hausse de l'offre
    \item \textbf{Progrès technique} : amélioration → hausse de l'offre
    \item \textbf{Nombre de producteurs} : plus de concurrents → hausse de l'offre
    \item \textbf{Réglementation} : contraintes → baisse de l'offre
    \item \textbf{Anticipations des producteurs} : anticipation de hausse → stockage
\end{itemize}

\section{Types de marchés}

\subsection{Marché de concurrence parfaite}
\pointcle{Marché théorique où le prix se forme librement}
\begin{itemize}
    \item Atomicité : nombreux acheteurs et vendeurs
    \item Homogénéité : produits identiques
    \item Transparence : information parfaite
    \item Libre entrée/sortie du marché
\end{itemize}

\subsection{Marchés imparfaits}
\pointcle{Marchés réels où la concurrence n'est pas parfaite}
\begin{itemize}
    \item \textbf{Monopole} : un seul vendeur
    \item \textbf{Oligopole} : quelques vendeurs dominants
    \item \textbf{Concurrence monopolistique} : produits différenciés
\end{itemize}

\section{Exemples concrets}

\exemple{Marché de l'essence : oligopole avec quelques grandes compagnies, prix similaires, demande peu élastique}

\exemple{Marché des smartphones : concurrence monopolistique, différenciation par marque, innovation constante}

\exemple{Marché agricole : proche de la concurrence parfaite, prix volatils selon offre/demande}

\section{Points clés à retenir}

\pointcle{Le prix est le signal qui guide les décisions économiques}

\pointcle{L'autorégulation du marché fonctionne par ajustement automatique des prix}

\pointcle{Les déséquilibres sont temporaires en concurrence parfaite}

\pointcle{L'élasticité détermine l'ampleur des variations de quantité selon les prix}

\pointcle{Les marchés réels s'éloignent souvent du modèle de concurrence parfaite}

\section{Mots-clés à connaître}
\begin{itemize}
    \item Marché
    \item Offre/Demande
    \item Prix d'équilibre
    \item Excédent/Pénurie
    \item Élasticité-prix
    \item Concurrence parfaite
    \item Monopole/Oligopole
    \item Biens substituts/complémentaires
    \item Autorégulation
    \item Mécanisme des prix
\end{itemize}

\end{document}